\vspace{-1em}
\section{Introduction}
\label{sec:intro}
Inverse rendering of glossy objects poses formidable challenges in computer vision and graphics. These surfaces exhibit strong view-dependent effects, such as specular highlights, environmental reflections, and inter-reflections, that inherently undermine multi-view consistency. This fundamental conflict severely impedes the accurate recovery of geometry and material properties from multi-view imagery. Although methods~\cite{mildenhall2021nerf, wang2021neus, barron2021mipnerf, barron2022mipnerf360, muller2022instantngp, chen2022tensorf, verbin2022refnerf, liu2023nero, liang2023envidr, jin2023tensoir} based on Neural Radiance Fields (NeRF) demonstrate promise in modeling complex light transport, their reliance on dense ray sampling compromises computational efficiency for practical deployment.

Recent advances in 3D Gaussian Splatting (3DGS)~\cite{kerbl20233dgs} deliver efficient scene reconstruction through explicit representations and real-time rendering. However, subsequent methods exhibit various limitations: certain approaches~\cite{guedon2024sugar, yu2024mipsplatting, yu2024gof, chen2024pgsr, lu2024scaffold, huang20242dgs} exhibit limitations for glossy objects, while others achieve plausible geometry reconstruction for glossy surfaces, yet either fail to perform material decomposition~\cite{tang20243igs, yang2024specgs, bi2024gs3, zhang2025refgs} or persistently yield blurred decomposition results despite enabling inverse rendering~\cite{jiang2024gaussianshader, liang2024gsir, gao2024r3dg, ye20243dgsdr, yao2024refgaussian, gu2025irgs}.
Specifically, geometry reconstruction is severely compromised by view-dependent specular reflections that violate multi-view consistency. This manifests itself as both high-frequency surface noise and structural artifacts, particularly in inter-reflection regions. Simultaneously, material decomposition is impaired by a simplified rendering equation and neglected indirect illumination, leading to blurred albedo, metallic, and roughness estimates that yield physically implausible relighting results.

We propose GOGS, a novel geometry prior-guided framework based on 2D Gaussian surfels that resolves these limitations through a two-stage pipeline. First, we establish robust surface reconstruction through physics-based rendering with split-sum approximation, enhanced by geometric priors from foundation models to explicitly enforce curvature continuity and mitigate ambiguities under specular interference. Subsequently, we perform physically accurate material decomposition on this optimized geometry by evaluating the full rendering equation via Monte Carlo importance sampling, where visibility and indirect radiance are computed through differentiable 2D Gaussian ray tracing. To address fidelity limit in specular rendering, we further introduce a spherical mipmap-based specular compensation mechanism that adaptively refines high-frequency details.

Extensive experiments demonstrate state-of-the-art performance in both geometry accuracy and material decomposition. Our contributions are as follows:
\begin{itemize}
    \item A robust geometric reconstruction method for glossy objects that mitigates geometry ambiguities by leveraging geometric priors and split-sum approximation; 
    \item Physically-based material decomposition via the full rendering equation evaluation with Monte Carlo importance sampling;
    \item An adaptive specular compensation mechanism that directionally refines high-frequency details, mitigating fidelity limits in specular rendering.
\end{itemize}